\honest{Zombie Responses}{Eliezer Yudkowsky}{2008}

Tired from yesterday's long post on zombies, I only reply to Richard Chappell's comment and fix a loose end.

Chappell says that what I call conceivable, others would call only ``apparently conceivable.'' I agree, and I want two different words, because the gap between ``I don't see a contradiction yet'' and ``this is logically possible'' is huge, NP-complete even in simple-seeming cases. Maybe ``apparently conceivable'' for what you get by closing your eyes and imagining zombies, and ``logically possible'' for what is established by exhibiting a complete model or proof; the argument for epiphenomenalism needs the second. \nb{Chalmers had the two words, prima facie and ideal conceivability, and the zombie argument as Chalmers states it uses the stricter one.}

Chappell calls my view type-A materialism. I decline the bundle until I agree with each part. ``What is consciousness?'' is a legitimate demand for insight, though the answer may partly be insights that show we asked it the wrong way. That is not eliminativism. It is realism about a problem that seems to need some solution, seems unable to have any, and is discussed in a way that depends heavily on the ad hoc architecture of human cognition.

Chappell says I showed the zombie world strange, not contradictory. I spell it out in seven steps. The zombie world, by definition, contains every cause of anything observable, including the cause of my saying ``I think therefore I am.'' When I focus my awareness on my awareness, my inner narrative shortly says so, and I can say it aloud. It seems that my awareness causes the narrative and the narrative moves my lips. The word ``consciousness,'' ``if it has any meaning at all,'' refers to what makes me say I have inward awareness. So the zombie world contains consciousness. But by definition it does not. And the third step seems very likely to be empirically true. My conclusion is ``a high empirical probability that the zombie world is logically impossible.'' You can save the zombie world only by letting something other than consciousness cause my talk of it, and that, combined with saying consciousness exists, is what struck me as deranged. Imagining that the word ``consciousness'' could have referred to some epiphenomenon, as ``water'' on Twin Earth refers to XYZ, does not make the zombie world possible. The zombie world must lack consciousness itself, not merely ``consciousness'': it must be a world without H2O, not a world without ``water.'' I hold it an empirical fact, given what ``consciousness'' refers to, that removing consciousness without moving atoms is logically impossible. \nb{Chappell replied in the comments that a premise of this kind is ``question-begging'' against the epiphenomenalist, and that ``it can't be an empirical question what's logically possible.'' In Chalmers's classification, resting the impossibility on an empirical fact is type-B materialism, not the type-A view the post declined, and Chalmers's two-dimensional argument was aimed at exactly the water analogy. The post does not address it.}

Chappell says a natural law guarantees the fit between experience and talk, so it is no miracle. But the law itself is the extra, improbable postulate: one posits a conscious inner world, an outer world that talks about consciousness for no reason, and that the two align perfectly, and the third does not follow from the first two. ``Miraculous'' was the wrong word in a philosophical context, where it means violating a law; I meant improbability from no source, as in belief in perpetual motion. Still, intuitively it is a miracle, on the order of a cracker taking on the substance of Christ's flesh while looking and behaving exactly like a cracker, guaranteed by a natural law. \nb{Chalmers grants the worry: on epiphenomenalism the fit ``seems to be something of a lucky coincidence.'' Chappell's comment went on to compare it with a universe fit for life, and the post does not answer that.}

Chappell says a zombie's words are meaningless. My answer: you can build an AI that refines an inner part of itself to correlate with its environment, with numbers that obey Cox's theorems and so deserve to be called probabilities; ``The Simple Truth'' gives my view of whether such an AI has beliefs. One AI outputs a map of Texas, another flies to Texas and chirps ``True'' or ``False'' as the highways match or not, all in physics. You may refuse to call it a map, but the AIs still chirp, and they would chirp ``False'' at Chalmers's ``belief in an epiphenomenal inner core.'' Outer Chalmers does all the mapping of reality; Inner Chalmers can only bless it with epiphenomenal meaning, unrelated to map-territory correspondence. So accuracy should be judged by looking at Outer Chalmers. \nb{The Texas check compares a map with something the checker can observe, and an epiphenomenal property cannot be observed. Chappell's comment ended with two points the post does not quote: the argument from reflective coherence counts against Chappell's view only if one assumes ``a causal theory of knowledge,'' and Chalmers does not commit to epiphenomenalism. Chalmers said the same on the blog three days later; Yudkowsky replied that the two theses entail each other.}

Then I fix yesterday's argument, which left out an assumption: a good AI enforces global rationality by enforcing local rationality, treating as a bug any part that, in its context, systematically adds beliefs that are not true when its condition holds. A causally closed system that believes in an epiphenomenal self would believe it in a zombie world too, so it is globally unreliable; a mind whose parts are all locally reliable is globally reliable; so some step in forming that belief must be locally unreliable, and self-inspection must find it. Without the local requirement, reflective coherence is too cheap: an AI that finds a part computing 2 + 2 = 5 while counting sheep would reason that 2 + 2 does seem to equal 5, so it had better keep the part. The general lesson: show that an argument is globally reliable because each step is locally reliable, not by comparing its conclusions with your intuitions. (In 2013 I added that this has since been formalized in ``Tiling Agents for Self-Modifying AI.'')

Last, a commenter tells me that N'Shama is related to the word for ``breath,'' not ``hear.'' Now that is a miraculously misleading coincidence: the word arose for different reasons but sounds exactly right to make me think of an inner listener. ``Oops.'' \nb{The commenter is right. The current text of ``Zombies! Zombies?'' still calls N'Shama ``the hearer.''}
